\documentclass[11pt]{article}

% ======================================================================
%  Paquetes (todos estándar en Overleaf)
% ======================================================================
\usepackage[utf8]{inputenc}
\usepackage[T1]{fontenc}
\usepackage[spanish,es-noquoting]{babel}
\usepackage{amsmath,amssymb,amsthm}
\usepackage{mathtools}
\usepackage[margin=2.6cm]{geometry}
\usepackage{tikz}
\usetikzlibrary{calc}
\usepackage{xcolor}
\usepackage{enumitem}
\usepackage[colorlinks=true,linkcolor=blue!55!black,citecolor=blue!55!black,
            urlcolor=blue!55!black]{hyperref}

% ----- entornos de teorema -----
\theoremstyle{plain}
\newtheorem{proposition}{Proposición}[section]
\newtheorem{corollary}[proposition]{Corolario}
\theoremstyle{definition}
\newtheorem{remark}[proposition]{Observación}

% ----- macros -----
\newcommand{\R}{\mathbb{R}}
\newcommand{\spanop}{\operatorname{span}}
\newcommand{\Msg}{\mathrm{MSG}}
\newcommand{\Agg}{\mathrm{AGG}}
\newcommand{\softmax}{\mathrm{softmax}}

\title{\textbf{Sobre la redundancia expresiva de la atención\\
sobre representaciones tipo NBFNet en \emph{link prediction}\\ de un solo grafo de conocimiento}}
\author{}
\date{Junio de 2026}

\begin{document}
\maketitle

\begin{abstract}
NBFNet es hoy el método de referencia para predicción de enlaces (\emph{link prediction})
en un único grafo de conocimiento. Una familia reciente de modelos le superpone una capa de
\emph{atención} (en el espíritu de los \emph{graph transformers}) con la expectativa de que aporte
una capacidad de razonamiento que el \emph{message passing} no posee. Esta nota responde, de forma
fundamentada, dos preguntas que delimitan el alcance de esa idea. \textbf{(P1)} ¿Un transformer cuyos
\emph{valores} provienen de un \emph{stream} tipo NBFNet captura algo que ese stream no contenga ya?
Mostramos que \textbf{no}: la atención es un promedio ponderado de sus valores y, por lo tanto, no
puede producir información en ninguna dirección ausente de todos ellos; como las representaciones de
NBFNet ya son la agregación óptima de la evidencia de caminos dentro de su horizonte, la atención solo
puede re-ponderar lo conocido. \textbf{(P2)} ¿Un transformer \emph{por sí solo}, sin NBFNet, puede
capturar lo que NBFNet captura? La respuesta es matizada: un transformer ciego a la estructura del
grafo no puede; un \emph{graph} transformer que use la adyacencia sí puede en principio ---pero al
hacerlo \emph{reimplementa} el message passing, y su flexibilidad adicional no se traduce en una
mejora medible en este régimen. Concluimos por qué, en estos datos, ``NBFNet ya captura todo lo que la
atención podría aportar''.
\end{abstract}

% ======================================================================
\section{Planteamiento}
% ======================================================================

Sea $G=(\mathcal{V},\mathcal{E})$ un grafo de conocimiento, con nodos (entidades) $\mathcal{V}$ y
aristas tipadas $(x,r,y)\in\mathcal{E}$, donde $r$ es un \emph{tipo de relación}. La tarea de
\emph{link prediction} consiste en, dada una consulta $(u,q,?)$ ---cabeza $u$ y relación $q$---,
puntuar cada candidato $v$ según la plausibilidad de la tripleta $(u,q,v)$.

NBFNet \cite{nbfnet} resuelve esta tarea con \emph{message passing} condicionado a la fuente, y es el
estado del arte en el régimen de \emph{un solo grafo} (transductivo, e inductivo sobre entidades).
Una propuesta natural ---y la que motiva esta nota--- es añadir una capa de \emph{atención} sobre las
representaciones de nodo que NBFNet produce, con la idea de que la atención permita comparar
candidatos entre sí, capturar dependencias de largo alcance, o componer relaciones de forma explícita.
Empíricamente, en \texttt{WN18RR} y \texttt{FB15k-237}, esa atención no mejora la métrica de forma
apreciable. La pregunta que abordamos es \emph{por qué}, y si hay una razón estructural.

Fijamos primero los dos objetos: NBFNet (Sección~\ref{sec:nbfnet}) y la atención
(Sección~\ref{sec:attn}).

% ----------------------------------------------------------------------
\subsection{NBFNet como agregación sobre caminos}
\label{sec:nbfnet}

NBFNet calcula, para cada candidato $v$, una representación $h_{u\to v}\in\R^{d}$ condicionada en la
fuente $u$ y la relación $q$, mediante una recursión de Bellman--Ford generalizado:
\begin{equation}
\label{eq:bf}
h^{(0)}_{u\to v} = \mathbf{1}[\,v=u\,]\,\boldsymbol{q},
\qquad
h^{(t)}_{u\to v} = \Agg\!\Big(\big\{\,\Msg\big(h^{(t-1)}_{u\to x},\, w_{r}\big)
\;:\; (x,r,v)\in\mathcal{E}\,\big\}\Big),
\end{equation}
con $\boldsymbol q$ el embedding de la consulta, $w_r$ los parámetros del tipo de relación $r$, y
$\Agg$ una agregación permutación-invariante. La inicialización $h^{(0)}_{u\to v}=\mathbf{1}[v=u]\,\boldsymbol q$
(el \emph{labeling trick}) ancla la propagación a la fuente.

El resultado central de \cite{nbfnet} es que, tras $L$ iteraciones, \eqref{eq:bf} computa una
\emph{suma generalizada sobre todos los caminos} de longitud $\le L$ entre $u$ y $v$:
\begin{equation}
\label{eq:pathsum}
h^{(L)}_{u\to v}
\;=\;
\bigoplus_{P\,\in\,\mathcal{P}^{\le L}_{u\to v}}\;
\bigotimes_{(x,r,y)\,\in\,P} w_{r},
\end{equation}
donde $\mathcal{P}^{\le L}_{u\to v}$ son los caminos $u\to v$ de a lo sumo $L$ aristas y
$(\oplus,\otimes)$ forman un semianillo. Dicho en palabras: \emph{el vector $h^{(L)}_{u\to v}$ ya es la
evidencia acumulada de todos los caminos relacionales de hasta $L$ pasos que conectan la cabeza con el
candidato.} Escribiremos $V_v := h^{(L)}_{u\to v}$ y lo llamaremos el \emph{valor} del nodo $v$. El
predictor final del baseline es \emph{puntual}: un lector $\rho(V_v)$ asigna un puntaje a cada $v$ por
separado.

Una propiedad que usaremos: como \eqref{eq:bf} es message passing de $L$ capas, el campo receptivo de
$V_v$ es la bola de radio $L$ alrededor de $\{u,v\}$. Por lo tanto $V_v$ solo puede contener evidencia
de caminos de longitud $\le L$; ningún camino más largo está representado en él
(Figura~\ref{fig:horizon}).

% ----------------------------------------------------------------------
\subsection{La operación de atención}
\label{sec:attn}

Una capa de atención actualiza la representación de cada nodo $v$ mezclando los valores de los demás:
\begin{equation}
\label{eq:attn}
o_v \;=\; \sum_{w} \alpha_{vw}\, u_w,
\qquad
\alpha_{vw} = \softmax_w\!\Big(\tfrac{1}{\sqrt{d}}\, q_v^{\!\top} k_w\Big),
\qquad
u_w = W_V\, V_w .
\end{equation}
Los \emph{queries} $q_v$ y \emph{keys} $k_w$ determinan los pesos, y $W_V$ proyecta los valores. Por
construcción del \emph{softmax}, los pesos son no negativos y suman uno:
\begin{equation}
\label{eq:weights}
\alpha_{vw}\ge 0 \quad\forall w,
\qquad
\sum_w \alpha_{vw} = 1 .
\end{equation}
La salida $o_v$ pasa luego por una conexión residual y una red \emph{feed-forward} puntual. La
propiedad \eqref{eq:weights} ---que la atención es un \emph{promedio ponderado} de los valores--- es
el núcleo del argumento de la Sección~\ref{sec:q1}.

\begin{figure}[t]
\centering
\begin{tikzpicture}[scale=1.0,>=stealth,
    nodo/.style={circle,draw,fill=blue!8,minimum size=15pt,inner sep=0pt,font=\small},
    fuente/.style={circle,draw,thick,fill=blue!25,minimum size=17pt,inner sep=0pt,font=\small},
    lejos/.style={circle,draw,dashed,fill=red!8,minimum size=15pt,inner sep=0pt,font=\small}]
  % horizonte L=2 (círculo punteado)
  \draw[gray!55,dashed] (0,0) circle (2.45cm);
  \node[gray!70] at (1.75,1.95) {\footnotesize horizonte $L$};
  % fuente
  \node[fuente] (u) at (0,0) {$u$};
  % nodos a 1 y 2 hops (dentro)
  \node[nodo] (a) at (1.3,0.5) {};
  \node[nodo] (b) at (0.6,-1.2) {};
  \node[nodo] (c) at (-1.4,0.3) {};
  \node[nodo] (d) at (2.2,-0.3) {$v_1$};
  % candidato lejano (fuera del horizonte)
  \node[lejos] (v) at (4.4,0.2) {$v_2$};
  \node[nodo] (e) at (3.3,-0.6) {};
  % aristas
  \draw (u)--(a); \draw (u)--(b); \draw (u)--(c); \draw (a)--(d);
  \draw (d)--(e); \draw (e)--(v); \draw (b)--(e);
  % etiquetas
  \node[blue!45!black,below] at (d.south) {\footnotesize dentro ($d(u,v_1)\le L$)};
  \node[red!55!black,below] at (v.south) {\footnotesize fuera ($d(u,v_2)>L$)};
\end{tikzpicture}
\caption{Campo receptivo de un NBFNet de $L$ capas. El valor $V_{v_1}$ de un candidato \emph{dentro}
del horizonte ($d(u,v_1)\le L$) contiene la evidencia de los caminos que lo conectan con la fuente; el
de un candidato \emph{fuera} ($d(u,v_2)>L$) no, porque la propagación de $L$ pasos no alcanza el camino
correspondiente.}
\label{fig:horizon}
\end{figure}

% ======================================================================
\section{P1: un transformer sobre valores de NBFNet es redundante}
\label{sec:q1}
% ======================================================================

Esta es, en esencia, la arquitectura cuestionada: NBFNet produce los valores $V_v$ y una capa de
atención los mezcla antes del lector. Mostramos que la atención no puede aportar información nueva, y
luego que, donde podría aportarla, esa información tampoco está en los valores.

% ----------------------------------------------------------------------
\subsection{La atención no crea direcciones ausentes de sus valores}

\begin{proposition}
\label{prop:span}
Sea $o_v=\sum_w \alpha_{vw}\, u_w$ la salida de la agregación de atención \eqref{eq:attn}, con pesos
que cumplen \eqref{eq:weights}. Entonces $o_v\in\spanop\{u_w\}_w$. En particular, si una dirección
$\varphi\in\R^{d}$ es ortogonal a todos los valores, es decir $\varphi^{\!\top}u_w=0$ para todo $w$,
entonces $\varphi^{\!\top}o_v=0$ \emph{cualquiera que sea el patrón de atención} $\{\alpha_{vw}\}$.
\end{proposition}

\begin{proof}
$o_v$ es una combinación lineal de los $u_w$, luego pertenece a su \emph{span}. La segunda afirmación
sigue por linealidad del producto interno:
$\varphi^{\!\top}o_v=\sum_w \alpha_{vw}\,\varphi^{\!\top}u_w=\sum_w \alpha_{vw}\cdot 0=0.$
\end{proof}

La Proposición~\ref{prop:span} formaliza una limitación simple pero decisiva: \textbf{la atención solo
\emph{recombina} los valores que recibe.} Si una característica discriminativa ---por ejemplo, la
evidencia de un camino particular--- no está presente en ningún valor $V_w$, ninguna distribución de
pesos la hace aparecer en la salida. La proyección $W_V$, la conexión residual y la red
\emph{feed-forward} posteriores actúan \emph{por nodo} (puntualmente) y no cambian esta conclusión: la
única vía por la que la representación de $v$ adquiere información sobre otro nodo $w$ es la mezcla
\eqref{eq:attn}, gobernada por la Proposición~\ref{prop:span}.

% ----------------------------------------------------------------------
\subsection{Consecuencia al combinarla con valores de NBFNet}

Cada $V_w$ es, por \eqref{eq:pathsum}, el agregado de todos los caminos $u\to w$ de hasta $L$ hops.
Conviene separar dos regímenes según la distancia $d(u,v)$ entre la cabeza y el candidato.

\paragraph{Dentro del horizonte, $d(u,v)\le L$.}
La evidencia de los caminos hacia $v$ está \emph{íntegra} en un \emph{único} valor, $V_v$, y el lector
puntual $\rho(V_v)$ ya la ve. Mezclar $V_v$ con otros valores no agrega evidencia de caminos: solo
re-pondera lo que ya estaba disponible. Esto explica por qué, en la práctica, quitar o alterar la
capa de atención no cambia apreciablemente la métrica: dentro del horizonte, la atención es
estrictamente redundante con la representación que NBFNet ya entrega.

\paragraph{Fuera del horizonte, $d(u,v)>L$.}
Ahora \emph{ningún} valor $V_w$ contiene el camino de longitud $>L$ que conecta la fuente con el gold,
porque el message passing de $L$ capas no llegó a propagarlo (Figura~\ref{fig:horizon}). Sea
$\varphi_{\text{largo}}$ la dirección que codificaría esa evidencia de largo alcance: por construcción
$\varphi_{\text{largo}}^{\!\top}V_w=0$ para todo $w$, y la Proposición~\ref{prop:span} implica que la
atención \emph{no puede recuperarla}. Es decir: justo donde la información faltante podría ayudar, no
está en los valores, y por lo tanto la atención sobre esos valores tampoco la puede explotar.

\begin{corollary}[La redundancia es una coincidencia de horizonte]
\label{cor:horizon}
Un transformer cuyos valores provienen de un NBFNet de horizonte $L$, y que opera \emph{dentro de ese
mismo horizonte}, es redundante: su salida es un promedio de agregados de caminos de longitud $\le L$
y, por la Proposición~\ref{prop:span}, no contiene evidencia de caminos ausente de todos los valores.
La única forma de aportar algo no redundante es \emph{introducir información de fuera del horizonte}
(caminos $>L$, o estructura por pares no codificada en los nodos). Pero esa información, por
definición, no es ``el valor que viene del stream NBFNet'': es un ingrediente adicional.
\end{corollary}

El Corolario~\ref{cor:horizon} también predice la \emph{única} forma en que la información relacional
de largo alcance puede ayudar: introduciéndola \emph{explícitamente} (por ejemplo, features de caminos
de longitud mayor que el horizonte del message passing), no recombinando los valores. Esto coincide con
lo que observamos: en las consultas cuyo candidato verdadero queda más lejos que el horizonte de
propagación ---donde la representación de nodo pierde poder discriminativo--- un descriptor explícito de
caminos largos aporta una señal complementaria, pequeña pero real.

\begin{remark}[Una objeción y su respuesta]
Podría argumentarse: ``si \emph{apilo} capas de atención \emph{sobre el grafo}, cada capa añade un hop
y así extiendo el horizonte''. Es cierto, pero entonces ya no se está mezclando un $V$ fijo de NBFNet:
se está haciendo \emph{más} message passing, ahora con pesos de atención. Eso corresponde a la
Pregunta~2 (Sección~\ref{sec:q2}, caso (b)) y, empíricamente, equivale a un NBFNet más profundo: no
recupera el régimen de largo alcance y tiende a degradarse por \emph{over-smoothing}
\cite{oono} (al apilar capas, las representaciones de los nodos se vuelven mutuamente
indistinguibles).
\end{remark}

% ======================================================================
\section{P2: ¿un transformer por sí solo captura lo que NBFNet captura?}
\label{sec:q2}
% ======================================================================

``Transformer por sí solo'' admite dos lecturas, y el matiz \emph{es} la respuesta.

% ----------------------------------------------------------------------
\subsection{Caso (a): transformer ciego a la estructura --- \emph{no}}

Si la atención es \emph{all-pairs} sobre embeddings de nodo \emph{sin} usar la adyacencia (los nodos
como un conjunto sin estructura), el modelo no tiene acceso a las aristas. Pero las aristas son lo que
\emph{define} los caminos, y los caminos son el contenido del razonamiento (ecuación
\eqref{eq:pathsum}). El modelo debería \emph{inferir} la adyacencia a partir de las features, lo cual
es inviable para grafos arbitrarios y, de forma más fundamental, \textbf{rompe la inductividad}: con
entidades nuevas en test (grafo disjunto), no hay embeddings de nodo aprendidos que transferir. Un
transformer ciego a la estructura no puede, por tanto, capturar NBFNet.

% ----------------------------------------------------------------------
\subsection{Caso (b): \emph{graph} transformer --- \emph{sí en principio, pero}}

Si la atención está \emph{restringida o sesgada por la adyacencia relacional} (cada nodo atiende a sus
vecinos, con un sesgo dependiente del tipo de relación), con la fuente marcada y $L$ capas apiladas,
entonces el modelo \textbf{puede igualar o exceder a NBFNet}. La razón es estructural:
\begin{itemize}[leftmargin=1.4em]
  \item \textbf{El message passing es un caso particular de la atención sobre aristas.} Un GNN es una
        capa de atención cuyo patrón está \emph{fijado} a la adyacencia; \emph{graph attention}
        \cite{gat} lo generaliza a pesos aprendidos. Así, $\{\text{message passing}\}\subseteq
        \{\text{atención restringida a aristas}\}$: una capa de \emph{graph attention} equivale a un
        hop, y $L$ capas a $L$ hops, con agregación aprendida en lugar de fija.
  \item \textbf{La composición relacional es expresable como atención.} El \emph{Edge Transformer}
        \cite{edgetransformer} usa atención \emph{triangular} sobre pares de entidades ---actualiza el
        par $(i,j)$ atendiendo a $k$ vía $(i,k)$ y $(k,j)$--- y compone relaciones $i\!\to\!k\!\to\!j$,
        igualando o superando a los GNNs en generalización sistemática. Un transformer \emph{puede},
        pues, realizar el razonamiento de caminos.
\end{itemize}

Hay, sin embargo, dos reparos decisivos.

\paragraph{(i) Para igualar a NBFNet hay que darle la estructura $\Rightarrow$ reimplementa el message
passing.} En cuanto el graph transformer usa las aristas, los tipos de relación y el nodo fuente, está
\emph{haciendo} message passing por medio de la atención. No lo trasciende: es ``NBFNet con agregación
aprendida en vez de fija''. La pregunta deja de ser ``¿puede?'' y pasa a ser ``¿la agregación
\emph{blanda y aprendida} aporta algo sobre la \emph{fija}?''.

\paragraph{(ii) Poder $\neq$ mejorar.} La versión expresiva (atención por pares / triangular) cuesta
$O(N^2)$ o $O(N^3)$ en el número de entidades $N$, prohibitivo en grafos grandes, lo que obliga a
atención \emph{sparse} o \emph{lineal}, que sacrifica justamente la expresividad extra. Y aun
teniéndola, la flexibilidad adicional resulta empíricamente inerte en \texttt{WN18RR} y
\texttt{FB15k-237}: la agregación fija de caminos de NBFNet ya es casi óptima para la masa de consultas
alcanzable, y el único margen restante (las consultas de largo alcance) no lo cierra la atención sobre
el mismo sustrato (Sección~\ref{sec:q1}).

% ======================================================================
\section{Síntesis y conclusión}
% ======================================================================

Los tres argumentos clásicos a favor de la atención sobre un GNN son: \textbf{(i)} alcance global en
una sola capa; \textbf{(ii)} agregación blanda dependiente del input, en vez de fija; \textbf{(iii)}
razonamiento por pares o de orden superior. En \emph{link prediction} de un solo grafo, los tres se
diluyen:
\begin{itemize}[leftmargin=1.4em]
  \item (ii) es \textbf{inerte}: promediar valores que ya son agregados de caminos no añade nada dentro
        del horizonte (Proposición~\ref{prop:span}); quitar la atención no cambia la métrica.
  \item (iii) es \textbf{redundante}: re-codificar pares o composición de relaciones reproduce ---peor
        codificada--- la información que el valor de nodo ya contiene.
  \item (i) es el \textbf{único margen real}, en las consultas de largo alcance, pero exige introducir
        información de fuera del horizonte de propagación (Corolario~\ref{cor:horizon}), no recombinar
        los valores; y en estos datos el efecto es pequeño.
\end{itemize}

\noindent
La conclusión no es que la atención sea débil, sino más sutil:

\begin{center}
\fbox{\parbox{0.94\linewidth}{\centering
NBFNet captura lo que la atención podría aportar porque su agregación fija de caminos ya es (casi)
óptima para la masa alcanzable: \emph{promediarla} no agrega información (P1), y \emph{reproducirla} con
atención es reimplementar el message passing (P2). La atención solo aportaría introduciendo evidencia
que el message passing estructuralmente no posee ---caminos más largos que su horizonte, o relaciones
por pares---, y ese margen, en estos benchmarks, es marginal.}}
\end{center}

\noindent
En términos de diseño: una capa de atención \emph{sobre} las representaciones de NBFNet es redundante
por construcción; un graph transformer \emph{en lugar} de NBFNet reimplementa el message passing sin
una ganancia medible; y la única dirección potencialmente fértil ---inyectar evidencia de caminos de
largo alcance, más allá del horizonte de propagación--- es, en \texttt{WN18RR}/\texttt{FB15k-237}, de
efecto pequeño. Esto acota con un argumento, y no solo con experimentos, qué puede esperarse de la
atención en este problema.

% ======================================================================
\begin{thebibliography}{9}
% ======================================================================
\bibitem{nbfnet}
Z.~Zhu, Z.~Zhang, L.-P.~Xhonneux, J.~Tang.
\emph{Neural Bellman-Ford Networks: A General Graph Neural Network Framework for Link Prediction}.
NeurIPS 2021. arXiv:2106.06935.

\bibitem{edgetransformer}
L.~Bergen, T.~O'Donnell, D.~Bahdanau.
\emph{Systematic Generalization with Edge Transformers}.
NeurIPS 2021. arXiv:2112.00578.

\bibitem{gat}
P.~Veli\v{c}kovi\'c, G.~Cucurull, A.~Casanova, A.~Romero, P.~Li\`o, Y.~Bengio.
\emph{Graph Attention Networks}. ICLR 2018. arXiv:1710.10903.

\bibitem{oono}
K.~Oono, T.~Suzuki.
\emph{Graph Neural Networks Exponentially Lose Expressive Power for Node Classification}.
ICLR 2020. arXiv:1905.10947.
\end{thebibliography}

\end{document}
